\chapter{First steps with \amlet{}}
\label{ch:launch}

\amlet{} can be used as a C library or a standalone application.
When used as a library, more complex models can be described and more analysis can be performed, but at the expense of some programming efforts.
The use as a standalone application is easier, but the range of models that can be estimated is more limited.

\section{Standalone application}

When \amlet{} is used as a command-line application, it has to be launched
from a text-console.
The basic usage is
\begin{verbatim}
amlet [OPTIONS] FILE
\end{verbatim}
where \verb|FILE| designs the driver file (see
Section~\ref{sec:driver}), and \verb|[OPTIONS]| corresponds to
optional options, as described in Table~\ref{tab:options}. The
options can be used in short or long form; for instance, the command
\begin{verbatim}
amlet -o log driver
\end{verbatim}
has the same effect that
\begin{verbatim}
amlet --output log driver
\end{verbatim}

\begin{table}[htb]
\begin{center}
\begin{tabular}{|c|c|c|c|}
\hline
short & long & argument & function \\
%\hline
\verb|-d| & \verb|--debug| & - & activate debug mode \\
\hline
\verb|-h| & \verb|--help| & - & print help message \\
\hline
\verb|-o| & \verb|--output| & file name
& save the log information into a file \\
\hline
\verb|-v| & \verb|--verbose| & - & verbose mode \\
\hline
\verb|-V| & \verb|--version| & - & print version number \\
\hline
\end{tabular}
\caption{command-line options}
\label{tab:options}
\end{center}
\end{table}

We briefly describe these options here, and cover them in more details
in the relevant sections.
\begin{description}
\item[debug] is meant for development purposes, delivering additional
  messages on the standard output in order to help debugging code;
\item[help] option makes \amlet{} displaying some brief help on the
  screen, and then exits.
\item[output] allows the user to define the name of the log
file (see Section~\ref{se:logfile}).
\item[verbose] option informs AMLET to print additional information on
  the screen during execution.
\item[version] prints the version number on the screen, afterwhat
the program exits.
\end{description}

\section{Driver file}
\label{sec:driver}

The \amlet{} driver file expresses a specification of the model to
calibrate, as well as of the optimization method to apply.
The driver has to be written in raw text format; examples of driver
files can be found in Appendix~\ref{app:driver}.
It basically consists of two kinds of lines: the command ones and the
affectations ones.
A command line has the following form:
\begin{verbatim}
command parameters;
\end{verbatim}
In other terms if consists of a command followed by the parameters to
apply to this command. An affection line has the form
\begin{verbatim}
variable = value;
\end{verbatim}
The value can be an integer, a real, a string or a set, depending on the
variable that has to be affected. In some cases, a predefined constant
identifier can be used instead of a numerical quantity. String values
are any sequence of characters, enclosed within double
quotes; a correct string would be for example
\begin{verbatim}
``this is a string''
\end{verbatim}
You cannot use double quotes inside a string.
We give in Table~\ref{tab:variables} the name of the main variables,
in alphabetical order. There will be discussed in more details in the
relevant sections.
\begin{table}[htbp]
\begin{small}
\begin{center}
\begin{tabular}{|c|c|p{6cm}|}
\hline
{\bf Parameter} & {\bf Default value} & {\bf Function} \\
\hline
\verb|b| & $(0,\ldots,0)^T$ & starting point \\
\hline
\verb|btype| & & type of the utilities coefficients \\
\hline
\verb|bgen| & & logit parameters for synthetic population generation \\
\hline
\verb|bgentype| & & type of the utilities coefficents during observations
generation\\
\hline
\verb|censor| & & 0 if all individuals face to all alternatives, 1
otherwise \\
\hline
\verb|iddep| & & dependent variable index\\
\hline
\verb|idcensor| & & variable index corresponding to the alternatives
availability \\
\hline
\verb|draws| & & indicates if new draws are to be generated or if they
have to be read from an external file\\
\hline
\verb|drawsname| & & name of the file containing the draws\\
\hline
\verb|eps| & & tolerance used for optimization \\
\hline
\verb|format| & \verb|Train96| & format of the data file \\
\hline
\verb|hessian_update| & & Hessian approximation method \\
\hline
\verb|level_accuracy| & & significance level used for objective accuracy
evaluation \\
\hline
\verb|method| & & optimization method \\
\hline
\verb|np| & & number of individuals \\
\hline
\verb|nalt| & & number of alternatives \\
\hline
\verb|niter| & & maximum number of iterations during the optimization
process \\
\hline
\verb|nobs| & & total number of observations \\
\hline
\verb|nrep| & & number of draws \\
\hline
\verb|nvar| & & number of variables (coefficients) \\
\hline
\verb|problem| & & problem name \\
\hline
\verb|robust| & & indicates if robust estimation of variance-covariance
matrix has to be used \\
\hline
\verb|rangen| & & random numbers generation method \\
\hline
\verb|sampling| & & sampling method \\
\hline
\verb|randomization| & & randomization strategy \\
\hline
\verb|seed1| & & seed of the random generator\\
\hline
\verb|validation| & \verb|false| & results validation \\
\hline
\end{tabular}
\caption{variables}
\label{tab:variables}
\end{center}
\end{small}
\end{table}

You can also write comments with the characters \# or \makeatletter
@\makeatother. The remaining of the line is then ignored by
\amlet{}, except if these characters are part of a string. The driver
file is case unsensitive, except for filenames, depending on the
operating system on which \amlet{} is executed.

\subsubsection{Commands}

We now briefly review the commands of \amlet{}. First, the command
\verb|output| expresses the name of file where the results will be saved:
\begin{verbatim}
output ``results'';
\end{verbatim}
where \verb|results| has to be replaced with the name that you want.
If no \verb|output| command is present in the drive file, the name
will be set by default to the name of the driver file followed by
\verb|.output|.

The next two commands allow to specify an existing observations file
or to generate a file that will be filled with synthetic data.
Assuming that your observations file has the name \verb|data|, you can
load it into \amlet{} with the command
\begin{verbatim}
load ``data'';
\end{verbatim}
On the other name, synthetic data can be obtained with the command
\begin{verbatim}
generate ``data'';
\end{verbatim}
The structure of the data files is explained in Section~\ref{sec:data}.
